% ---
% id: protocol-compiler/thm:validator-soundness
% kind: thm
% title: "Soundness, completeness and decidability of the validator"
% summary: >-
%   Passing the validator's eight checks establishes every target condition of
%   a well-formed declaration, and every check but the feasibility check is
%   complete, so a well-formed declaration is refused only when no default or
%   supplied candidate is verified feasible. Every decision is finite and exact.
% prominence: main
% status: implemented
% uses: [thm:feasibility, thm:measurements-and-statistics, def:declaration, def:validator, def:homog, def:target, def:exact, def:feasible]
% mentions: [rem:nosearch, thm:band-a, def:chaincond, thm:sifting-and-key, req:operator-form, foundations/req:exact-island]
% assumes: []
% refs: []
% provenance: newly-derived
% implements: certiqs_key.protocol_compiler.validation:decidability_counts
% lean:
% notes: "thm:sifting-and-key is a mention and not a direct use: the general-format Step 3 names its (v) as a pointer, and no step of the proof rests on it. def:chaincond and thm:band-a are mentions: the proof's pointer to (C1) says where the descriptor equality is decided, V1 reading no chain descriptor, and (iv)'s pointer to thm:band-a (i) is explanatory. The counts of (iv) are this theorem's own assertions about the checks it pairs, not restatements of a tool's output. Also realised by certiqs_key.protocol_compiler.validation:reconstruct_sifting_predicate and reconstruct_holder_key_rule."
% ---

The consuming engine takes D(i)--D(viii) as hypotheses without re-deriving
them, so what a pass of the validator buys has to be exactly those
conditions. The theorem pairs each check of Definition~\ref{def:validator}
with its target condition of Definition~\ref{def:target}, names the one check
that is not equivalent to its condition, and counts what the checks compute.
Remark~\ref{rem:nosearch} records why an empty \tref{feasible set} is the dangerous
error and a spurious refusal is not. For the clause of (iv) that the
validator assembles no array larger than one matrix of size $d_A$, $d_s$ or
$d_Ad_s$ for a single sector $s$, see also
Requirement~\ref{req:operator-form}, by which the engine holds every operator
on $\Qsys$ as its block list.

\begin{theorem}
\label{thm:validator-soundness}
For a candidate declaration with finite tabulated data:
\begin{enumerate}[label=(\roman*),leftmargin=2.2em]
\item \emph{(soundness, on all eight conditions)} if V1--V8 all
pass then D(i)--D(viii) all hold, of the declared numbers;
\item \emph{(completeness, on seven of the eight)} if D(i)--D(vi)
and D(viii) hold then V1--V6 and V8 all pass;
\item \emph{(V7 is one-sided, and the gap is exactly its third branch)} the
converse of (i) fails, and it fails only at V7: D(vii) does \emph{not} imply
V7. For a witness, take \tref{untrusted-source mode} with $\Qsys$ a qubit
carrying the single sector $\SecSet=\{s_0\}$, every other slot well-formed,
and the one-entry list
$\CList=\{(\dyad{0}_{\Qsys},\ \ge,\ 1)\}$, whose entry reads
$\tr[(\idon{A}\otimes\dyad{0})\,\omega]\ge1$. Then $\hat O_1$ is a
one-block list, Hermitian, the typing is that of D(vi), and
$\omega=\eta\otimes\dyad{0}$ is feasible for any state $\eta\in\Seq{A}$ --- it
is Hermitian and positive semidefinite outright, has unit trace, and its
constraint value is exactly $1$ --- so $\Feas\ne\varnothing$ and D(i)--D(viii)
all hold. But $\CList\ne\varnothing$, and the default candidate
\eqref{eq:witness-empty} that branch (b) of V7 tries first gives
$\tr[(\idon{A}\otimes\dyad{0})(\idon{A}/d_A\otimes\idon{\Qsys}/2)]=\tfrac12$,
which fails $\ge1$; so with no witness supplied branch (c) is unavailable, V7
fails and the \tref{well-formed declaration} is refused. Exactly the same
happens for every declaration with $\CList\ne\varnothing$ presented without a
witness whose \tref{feasible set} is non-empty but does not contain
\eqref{eq:witness-empty}. The refusal errs on the safe side --- a refused
declaration is never consumed --- and the incompleteness is deliberate: deciding D(vii) in general is a
semidefinite feasibility problem this module does not solve;
\item \emph{(decidability, and what is computed)} each of V1--V6 and V8 is
decided by finite enumeration over the declared alphabets together with
finitely many arithmetic operations on the declared matrices; V7 is decided in
the same sense on the candidate of its branch (b) or (c), and is otherwise a
refusal. The arithmetic part consists of \emph{exact decisions}: an
identity, a semidefiniteness, or the exactness of an addition. The following
counts give the complete applicable predicate inventory for well-typed data,
not the number an early-terminated execution necessarily performs.
Write $I_{\mathrm{tr}}=1$ in trusted-source mode and $I_{\mathrm{tr}}=0$ when
the marginal slot holds $\nomarg$. The complete arithmetic inventory of V2,
V3 and V6 is
\begin{equation}\label{eq:predcount}
  \underbrace{3I_{\mathrm{tr}}}_{\text{V2}}
  \;+\;\underbrace{2|\Xset|+1}_{\text{V3 on }A}
  \;+\;\underbrace{|\SecSet|\,\big(2|\Vset|+1\big)}_{\text{V3 on }\Qsys}
  \;+\;\underbrace{m\,|\SecSet|}_{\text{V3 on }\CList}
  \;+\;\underbrace{m_{\text{gen}}}_{\text{V6}}
\end{equation}
predicate decisions, where $m_{\text{gen}}\le m$ counts the entries declared
general --- the only ones carrying a matrix whose Hermiticity is
decided; V6's remaining content is $m$ comparisons of a relation symbol, a tag
and a type, which are decided and not computed. \emph{Beside these stand the
summation decisions of condition~(S)}: one on $\tr\tau_A$ in trusted
mode and none on the absent marginal in untrusted mode; one on
$\sum_xM_x$; and one per sector on $\sum_vF_v^{(s)}$, so
$|\SecSet|+1+I_{\mathrm{tr}}$ of them,
each a pass over the additions of one declared-order sum. And V7, on the
candidate it verifies, takes $2|\SecSet|+1$ decisions --- Hermiticity and
semidefiniteness per block, and the trace --- one further in trusted mode for
the marginal equality, and one decision per constraint entry: semidefiniteness
for an inequality, exact entrywise equality for the relation $=$
(Theorem~\ref{thm:feasibility}\,(ii)); in untrusted mode a constraint leg is one
comparison of real numbers instead. Every one of them is taken on a matrix of
size $d_A$, $d_s$ or $d_Ad_s$ for a single $s$; \emph{the validator
assembles no array larger than one such matrix}, which where $|\SecSet|=1$ is
$d_{\Qsys}$ or $d_Ad_{\Qsys}$ square, the one block being the whole space
there (Theorem~\ref{thm:band-a}\,(i)).
In trusted mode the V2 term is $3$ and not $4$: D(ii)'s named Hermiticity
condition and the Hermiticity half of its positivity condition are the same
decision, taken once. In untrusted mode V2 decides only the marginal slot's
$\nomarg$ branch, with no marginal arithmetic or marginal summation predicate.
\end{enumerate}

\emph{The general format.} All four parts hold for a candidate declaration in
the general format of Definition~\ref{def:declaration}\,(DX1)--(DX6), with the
checks read as in the general-format paragraph of
Definition~\ref{def:validator} and the conditions as in that of
Definition~\ref{def:target}: (i) soundness on all eight and (ii) completeness
on D(i)--D(vi) and D(viii); (iii) with the \emph{same} witness, which is a
general-format declaration at $|\mathbb{T}|=1$, one \tref{measured factor} and
no party-local entry; and (iv) with the predicate decisions of V2, V3 and V6
numbering
\begin{equation}\label{eq:predcount-general}
  \underbrace{I_{\mathrm{tr}}\big(2|\mathbb{T}|+1
    +|\mathbb{T}|\,\mathbb 1[\text{ensemble form}]\big)}_{\text{V2}}
  \;+\;\underbrace{|\mathbb{T}|\,\big(2|\Xset|+1\big)}_{\text{V3 on }A}
  \;+\;\underbrace{\sum_{f\in\mathbb{F}}|\SecSet^{f}|
    \big(2|\Vset^{f}|+1\big)}_{\text{V3 per measured factor}}
  \;+\;\underbrace{\sum_{k=1}^{m}|\SecSet^{f_k}|}_{\text{V3 on }\CList}
  \;+\;\underbrace{m_{\text{gen}}\,|\mathbb{T}|}_{\text{V6}} ,
\end{equation}
beside them the summation decisions of condition~(S) --- one on
$\tr\tau_A$ in trusted mode only, one per Alice sector on Alice's family,
and one per factor and sector on that factor's family, so
$|\mathbb{T}|+\sum_f|\SecSet^{f}|+I_{\mathrm{tr}}$ of them --- and V7, on the
candidate it verifies, $2|\mathbb{T}|\,|\SecSet|+1$ decisions, $|\mathbb{T}|$
further in trusted mode for the marginal equalities, and per entry $k$ either
$|\mathbb{T}|$ operator comparisons, one comparison of real numbers in
untrusted mode, or, for a party-local entry,
$|\mathbb{T}|\,|\SecSet|/|\SecSet^{f_k}|$ operator comparisons.
Each operator comparison is one semidefiniteness decision for an inequality
or one exact entrywise equality decision for $=$
(Theorem~\ref{thm:feasibility}\,(ii), general format). \emph{Every one of them is
taken on a matrix of size $d_{A_t}$, $d_s$ or $d_{A_t}d_s$ for a single joint
sector, or of size $d_{A_t}d_s/d_{f,s_f}$ for a single party-local row;
the validator assembles no array larger than one joint block.} In the
\tref{basic format} --- $|\mathbb{T}|=1$, one measured factor, POVM form, two
messages neither reading the other and no party-local entry ---
\eqref{eq:predcount-general} is
\eqref{eq:predcount} term for term, and every other count of this paragraph is
the corresponding count of (iv).
\end{theorem}

\begin{proof}
\emph{Idea.} Parts (i) and (ii) are proved together, by pairing each check
with its target condition. Three kinds of conjunct arise: combinatorial ones,
which both sides state as the same comparison of labels, integers or
tabulated values; numerical ones, which both sides state as the same exact
predicate on the same declared numbers; and structural ones, which the block
type makes true and neither side states. Two pairs, V5/D(v) and V8/D(viii),
are combinatorial but need an argument. V7 implies D(vii), which completes
(i), and (iii) is the counterexample stated with it. Part (iv) counts the
decisions. The general format takes the same steps, the point being that on
each side the quantifiers still run over one list.

\emph{Step 1 (combinatorial conjuncts: the restatement is literal).} V1's
finiteness, totality, dimension-sum, index-order and sector-map clauses; V4's
key-alphabet clause; V6's relation-symbol, tag and type clauses. Each compares
labels, integers or tabulated values, so V1/D(i), V4/D(iv) and the
combinatorial half of V6/D(vi) are equivalences in both directions. For V1's
sector-map clause we \emph{name} both lists rather than identify them: D(i)
asks that there be one map and that every declared $\Qsys$-operator be a
block list against \emph{it}, and the declared $\Qsys$-operators of
\eqref{eq:declaration} are exactly the family $\{F_v\}_{v\in\Vset}$ together
with the $m$ entries $\hat O_1,\dots,\hat O_m$ of $\CList$
(Definition~\ref{def:declaration}); V1 compares the four components of the
map carried by \emph{each of those same objects} against the emitted one, and
compares the map of nothing else (Definition~\ref{def:validator}). The two
quantifiers therefore run over one list, and two finite lists of labels and
integers agree or they do not. \emph{Neither the descriptor that the chain's
last link produces nor a supplied feasibility witness is on either side}: the
first is a
component of the chain record, decided at
Definition~\ref{def:chaincond}\,(C1), and the second is a block list on
$A\Qsys$ tested at branch~(c) of V7 --- so a candidate meeting D(i)--D(vi)
and D(viii) cannot be refused at V1 on account of either, which is what (ii)
asserts and what a wider V1 would have made false. \emph{V1's index-order
clause is the last clause of D(i)}, word for word: both ask that the
restriction of $\idxm$ to $\SecSet$ be the total order $\Smap$ fixes, and
that comparison is again of two finite lists of labels. Both directions
therefore hold of it as well, which is what makes (ii) true of a
\emph{candidate}: a candidate meeting D(i)--D(vi) and D(viii) meets that
clause too, so V1 cannot refuse it there.

\emph{V6's combinatorial clauses pair with D(vi) on the same $m$ objects.}
D(vi) asks of each entry of $\CList$ that its relation symbol lie in the set
its emitted source mode admits and that its right-hand-side slot hold one of
the two admitted data, whose type \emph{is} the tag
(Definition~\ref{def:homog}); V6 decides those two things, of those same $m$
entries, keyed on whether the marginal slot is filled --- which is the
emitted source mode (Definition~\ref{def:declaration}). V6's remaining
clause, that an explicit tag field on an \emph{arriving} entry agree with the
type of the datum it names, is not a further conjunct on
\eqref{eq:declaration}: no entry of a declaration carries such a field, so on
a candidate that clause has no subject and cannot refuse one
(Definition~\ref{def:validator}, V6). The pairing is therefore between equal
lists here as well.

\emph{Step 2 (numerical conjuncts: the two sides are the same exact
predicate).} V2's three, V3's five and V6's Hermiticity of a $\Theta_k$
declared general are, by construction, the same conditions on the same
declared numbers as the corresponding conjuncts of D(ii), D(iii) and D(vi):
both Definition~\ref{def:target} and Definition~\ref{def:validator} state
them through Definition~\ref{def:exact}, exactly, with no window on either
side. So V2/D(ii), the five-predicate half of V3/D(iii) and the Hermiticity
half of V6/D(vi) are equivalences in both directions, a literal restatement.
Condition~(S) pairs the same way: it is a clause of D(ii) and of D(iii)
(Definition~\ref{def:target}) and a clause of V2 and V3, stated once in
Definition~\ref{def:exact} and read by both.

\emph{Step 3 (structural conjuncts: there is nothing left to pair).} D(iii)'s
last two conjuncts, $\Pi_sF_v\Pi_{s'}=0$ and $\Pi_s\hat O_k\Pi_{s'}=0$, hold
by the emitted type (Theorem~\ref{thm:measurements-and-statistics}\,(ii)) and have no counterpart in
V3. They are therefore vacuous in the soundness direction --- V3 need not
establish what cannot fail --- and vacuous in the completeness direction ---
a declaration satisfying them is not thereby subjected to any check. The
pairing is not broken by their absence; it is one conjunct shorter on each
side. One direction of (ii) does fail for a reason that is not numerical, and
it is a \emph{typing} failure: an operator presented outside the block type,
as a full $d_{\Qsys}\times d_{\Qsys}$ array, is refused at V1 whatever its
off-block entries are, exactly zero included, because it is not a declaration
in the sense of Definition~\ref{def:declaration}. That is why (ii) is stated
for declarations rather than for operators.

\emph{Step 4 (V5 $\Leftrightarrow$ D(v) and V8 $\Leftrightarrow$ D(viii)).}
These two pairs are combinatorial but need an argument.

V5 $\Leftrightarrow$ D(v). ($\Leftarrow$) If $\varsigma$ exists as in
D(v) and $a(x,v)=a(x',v')=\kappa$, then $r(x,v)=\perp$ and $r(x',v')=\perp$ are
both equivalent to $\varsigma(\kappa)=\mathrm{fail}$, so V5 holds.
($\Rightarrow$) Assume V5 and define
$\varsigma(\kappa)\coloneqq\mathrm{fail}$ if there exists
$(x,v)\in a^{-1}(\kappa)$ with $r(x,v)=\perp$, and
$\varsigma(\kappa)\coloneqq\mathrm{pass}$ otherwise (in particular on
$\kappa\notin\operatorname{ran}a$, where the choice is immaterial because the
condition of D(v) quantifies over pairs $(x,v)$). This $\varsigma$ is total.
For $(x,v)$ with $\kappa=a(x,v)$: if $r(x,v)=\perp$ then
$\varsigma(\kappa)=\mathrm{fail}$ by construction; if
$\varsigma(\kappa)=\mathrm{fail}$ then some $(x',v')\in a^{-1}(\kappa)$ has
$r(x',v')=\perp$, and V5 gives $r(x,v)=\perp$. Hence D(v).

V8 $\Leftrightarrow$ D(viii). ($\Leftarrow$) If \eqref{eq:partycomp}
holds with a total $r_{\hold}$, then equal holder readouts and equal
announcements give equal key values, which is V8; the announcement clause of
D(viii) is the first clause of V8. ($\Rightarrow$) Assume V8 and define
$r_{\hold}(o,\kappa)\coloneqq r(x,v)$ for any $(x,v)$ with
$\rd{\hold}(x,v)=o$ and $a(x,v)=\kappa$ if such a pair exists --- well defined
precisely by V8's implication --- and $r_{\hold}(o,\kappa)\coloneqq\perp$ on
pairs $(o,\kappa)$ that are not realized. This $r_{\hold}$ is total on
$\Oset{\hold}\times\Kann$ and satisfies \eqref{eq:partycomp} for every $(x,v)$,
since every such pair realizes $(\rd{\hold}(x,v),a(x,v))$. Hence D(viii). The
value chosen on unrealized pairs is immaterial: \eqref{eq:partycomp} quantifies
over $(x,v)$ only.

Steps 1 to 4 account for (ii) in full, and for (i) except at V7.

\emph{Step 5 (V7 $\Rightarrow$ D(vii), which completes (i); and (iii)).}
Assume V1--V8 all pass. By the equivalences of Steps 1 to 3, D(i)--D(iii)
hold. If V7 passed by branch (a), $\CList=\varnothing$ and
Theorem~\ref{thm:feasibility}\,(i) applies --- its hypotheses are exactly
D(i)--D(iii) --- and exhibits $\omega^{\star}\in\Feas$. If V7 passed by
branch (b) or (c), the candidate it verified satisfies every condition of
Definition~\ref{def:feasible}, decided exactly (Theorem~\ref{thm:feasibility}\,(ii),
whose further hypothesis D(vi) holds because V6 passed, by Steps 1 and 2),
so that candidate lies in $\Feas$. Either way D(vii) holds. Statement (iii) is
proved by the counterexample given with it, which is checked there against
Definition~\ref{def:target} condition by condition.

\emph{Step 6 (decidability, (iv)).} V1 enumerates the declared tables and
compares finitely many declared integers and labels, the sector-map
comparison included: it is $|\SecSet|$ identifier comparisons, $|\SecSet|$
dimension comparisons, one block-kind comparison and one length comparison,
per supplied family, plus the comparison of $\idxm$'s sector order with
$\Smap$'s. V2, V3 and V6 take the decisions \eqref{eq:predcount} counts, on
matrices of the declared finite dimensions: in trusted mode $3$ for V2 ---
$\tau_A=\tau_A^{\dagger}$, $\tau_A\succeq0$, $\tr\tau_A=1$, the first serving
both the Hermiticity condition of D(ii) and the Hermiticity half of its
positivity condition; in untrusted mode none, because V2 takes its
$\nomarg$ branch; $2|\Xset|+1$ for V3's Alice half, one Hermiticity and one
semidefiniteness per $x$ and one completeness identity on $A$;
$|\SecSet|(2|\Vset|+1)$ for its Bob half, the same per $(v,s)$ and one
completeness identity per $s$; $m\,|\SecSet|$ Hermiticity decisions for its
constraint half; and $m_{\text{gen}}$ for V6. Each is finite: an identity is
a comparison of entries, and a semidefiniteness on a Hermitian matrix of size
$n$ is decided by exact symmetric elimination with the zero-pivot row test,
$O(n^3)$ exact rational operations. The summation decisions of (S) are one
pass over the additions of each applicable declared-order sum; the marginal
trace sum exists only in trusted mode. Early termination can omit later
decisions within a check, which is why this complete inventory need not
equal the length of a failed execution's transcript. V4, V5 and V8 enumerate
$|\Xset||\Vset|$ pairs, respectively $(|\Xset||\Vset|)^2$ ordered pairs, and
compare tabulated values. V7 is Theorem~\ref{thm:feasibility}\,(ii), which is likewise
per block, and \emph{on both of its branches the decision is taken in exact
arithmetic}: on branch (b) the candidate is formed by the compiler and decided
in exact rational arithmetic, and on branch (c) the supplied witness's
membership is decided exactly on the dyadic numbers it declares --- not
because those numbers are exactly representable, which by
Definition~\ref{def:exact} warrants nothing about the additions and products
a decision forms, but because every decision this module takes on a computed
quantity is taken exactly (\xref{foundations/req:exact-island}).

\emph{The general format.} The three kinds of conjunct are the same three, and
\emph{no fourth kind arises}; what has to be checked is that on each side the
quantifiers still run over one list. The steps are those above.

\emph{Step 1 in the general format (combinatorial conjuncts).} V1's
finiteness, totality, factor-list, dimension, index-order and map-comparison
clauses; V4's; V6's relation-symbol, tag, type and \emph{scope} clauses. Each
compares labels, integers or tabulated values, so each is an equivalence in
both directions once the two quantifiers are seen to run over one list, which
needs saying in four places. \emph{First, the factor list.} D(i) asks that
the sector map be one factorised map --- at most one factor tagged $A$ and
then first, at least one tagged $\Qsys$, one block kind, the dimension
identities --- and V1 decides those same clauses on the same list
(Definition~\ref{def:validator}, the general format), so they pair literally.
\emph{Second, the objects whose maps are compared.} In the general format the
declared $\Qsys$-operators are the per-factor families
$\{F^{f}_{v_f}\}_{v_f}$, one per $f\in\mathbb{F}$, together with the $m$
operators $\hat O_1,\dots,\hat O_m$, each against the map of its scope
(Definition~\ref{def:declaration}\,(DX1), (DX3), (DX6)); \emph{and Alice's
declared data are on the list beside them} --- $\tau_A$, every operator of her
measurement slot, and every $\Theta_k$ of an entry declared general, each
against the map of the factor tagged $A$, which is where
Definition~\ref{def:exact} states condition~(T) non-empty for them. D(i) asks
that there be one factorised map and that \emph{each} of those objects be a
block list against the map of its own factor; V1 compares the four
components of the map carried by each of \emph{those same} objects and of
nothing else. \emph{So the Alice-side conjunct is on both sides and pairs like
the others}, which is what makes (i) sound and (ii) complete for it: a
candidate whose Alice block list carries a different block kind, a different
identifier list or a different order is refused at V1, and a candidate
meeting D(i) is not refused there, both sides comparing two finite lists of
labels and integers. Where no factor is tagged $A$, Alice's data are single
arrays carrying no map and both sides are empty of them, as in Step 1 above.
\emph{A product block \eqref{eq:productfamily} and the lifted operator of a
scoped entry are on neither side}: neither is a \tref{declared datum}, so D(i) does
not quantify over either and V1 compares neither --- which is what keeps (ii)
true, a wider V1 refusing candidates D(i)--D(vi) and D(viii) admit. The
descriptors that the chain's last links produce, one per factor, and a
supplied witness stay outside both sides for the reasons of Step 1, now read
per factor and over the joint sectors. \emph{Third, the index order.} D(i) asks
that the restriction of $\idxm$ to each factor's sectors be the order that
factor's map fixes, and records that the product order is \emph{derived}; V1
decides exactly that and compares nothing about the derived orders, there
being no second datum to compare them against. Both directions hold, so a
candidate meeting D(i)--D(vi) and D(viii) cannot be refused there.
\emph{Fourth, the transcript.} D(viii) asks that the announcement slot be a
\tref{transcript protocol} with every table $a_m$ total on
$\Oset{\pi_m}\times\mathcal{C}_1\times\dots\times\mathcal{C}_{m-1}$ --- so that
a message reads its sender's readout and the earlier messages \emph{by its
signature} --- and V1's totality clause together with V8's announcement clause
decide that same thing of those same $L$ tables. A table whose stated domain
names a later message's alphabet is not of the declared type and is refused at
V1; so on a candidate the signature clause is not a numerical conjunct and
needs no pairing beyond this one.

\emph{And V6 pairs with D(vi) on the same $m$ entries, now with three forms and
a scope.} D(vi) asks of each entry that its relation symbol lie in the set
\emph{its form and its emitted mode} admit --- the mode's set for an untagged
entry, $\{\preceq,=,\succeq\}$ in either mode for an entry declared
party-local, with the scalar symbols also admitted where every row of that
entry is $1\times1$ --- that its right-hand-side slot hold one of the three
admitted data --- whose type, together with the arriving tag field for the
party-local form, \emph{is} the tag (Definition~\ref{def:homog}) --- and that
its scope name a measured factor of the map; V6 decides those three things of
those same $m$ entries, keyed on whether the marginal slot is filled and on
the declared tag, and its symbol clause is D(vi)'s symbol clause set for set,
so that pair is an equivalence in both directions as before. The clause of V6
about an arriving tag field is again not a further conjunct on
\eqref{eq:declaration} for a homogeneous or general entry, whose tag is the
type; \emph{for a party-local entry it is the tag, because the type of the
datum does not separate that form from the homogeneous one}
(Definition~\ref{def:homog}), and D(vi) in its general reading asks for that
declared form, so here too the two sides ask the same thing of the same
objects.

\emph{Step 2 in the general format (numerical conjuncts).} V2's, V3's and
V6's Hermiticity of a $\Theta^{(t)}_k$ declared general are, by construction,
the same exact predicates on the same declared numbers as the corresponding
conjuncts of D(ii), D(iii) and D(vi), both sides stating them through
Definition~\ref{def:exact} in its general-format reading, exactly and with no
window: per measured factor and per factor block on the measured families,
and per Alice sector on Alice's data. Condition~(S) pairs the same way, per
factor and per Alice sector. \emph{Two pairs are new and both are
equivalences.} The ensemble form: V3's identity $\sum_xN^{(t)}_x=\tau^{(t)}_A$
with its summation condition, and V2's $\tau^{(t)}_A\ne0$, are exactly the
conditions Definition~\ref{def:exact} states for that form and D(ii), D(iii)
ask; and V3's refusal of an ensemble-form slot beside a marginal slot holding
$\nomarg$ pairs with D(iii)'s clause that ensemble form is refused in
untrusted-source mode --- a decision on \emph{which of two slots is filled},
hence combinatorial, hence an equivalence in both directions.

\emph{Step 3 in the general format (structural conjuncts: again nothing is
left to pair, and the list of them is longer).} D(iii)'s typing conjuncts
hold by the emitted type along $\mathbb{T}$ as well as along $\SecSet$ and
have no counterpart in V3; and \emph{three further statements that a naive
reading would put on the data are theorems and are conjuncts of neither
side}: that the product family $\{F_v\}_{v\in\Vset}$ is a POVM on $\Qsys$,
block diagonal for the product map (Definition~\ref{def:declaration}\,(DX3));
that the family \eqref{eq:ensemblepovm} an ensemble form denotes is a POVM on
$A$ --- proved at Theorem~\ref{thm:sifting-and-key}\,(v), \emph{named here as a pointer
and used by no step below}, what this step needs being only that neither
D(iii) nor V3 states it; and that a lifted operator is Hermitian and block
diagonal for the product map, which follows from the Hermiticity of
$\hat O^{(s_{f_k})}_k$ that V3 does decide. So each is vacuous in the
soundness direction --- V3 need not establish what follows --- and vacuous in
the completeness direction, a declaration satisfying it being subjected to no
check for it. The typing failure that makes (ii) a statement about
declarations rather than operators is unchanged, and has two more instances: a
family presented as one array over the product alphabet, and an Alice operator
presented as one array on $A$ where an $A$ factor is declared, are refused at
V1 whatever their entries are, because neither is a declaration in the sense
of Definition~\ref{def:declaration}.

\emph{Steps 4 and 5 in the general format (V5, V8 and V7; and (iii)).} The
proofs of V5 $\Leftrightarrow$ D(v) and V8 $\Leftrightarrow$ D(viii) in
Step 4 use of $a$ only that it is a total map into $\Kann$, and construct
$\varsigma$ and $r_{\hold}$ on that alphabet; both hold verbatim with
$\Kann=\prod_m\mathcal{C}_m$ and $a$ the function \eqref{eq:transcript}
derives, \emph{whose totality is derived here and not imported}:
$\kappa_1=a_1\circ\rd{\pi_1}$ is total because $a_1$ and $\rd{\pi_1}$ are, and
if $\kappa_1,\dots,\kappa_{m-1}$ are total on $\Xset\times\Vset$ then so is
$\kappa_m$, $a_m$ being tabulated on all of
$\Oset{\pi_m}\times\mathcal{C}_1\times\dots\times\mathcal{C}_{m-1}$ (V1) and
$\rd{\pi_m}$ total; induction on $m$ gives $a$ total into $\Kann$, which is
all the two proofs consume of it. For V7 $\Rightarrow$ D(vii): if V7 passed by
branch~(a) then $\CList=\varnothing$ and Theorem~\ref{thm:feasibility}\,(i) in its
general-format reading exhibits \eqref{eq:witness-general} in $\Feas$, its
hypotheses being exactly D(i)--D(iii); if by branch~(b) or (c), the candidate
it verified satisfies every line of \eqref{eq:feas-general}, decided exactly
(Theorem~\ref{thm:feasibility}\,(ii), general format, D(vi) holding because
V6 passed), so it lies in $\Feas$. Either
way D(vii) holds, and (i) is complete. \emph{And (iii) holds in the general
format with the witness already given}: that declaration has one Alice sector,
one measured factor, POVM form, two messages and no party-local entry, at which
values D(i)--D(viii) in their general reading are the eight conditions read of
it in (iii) (Definition~\ref{def:target}, the general format) and V7's three
branches are the three branches read of it there, so it is a general-format
declaration satisfying D(i)--D(viii) that V7 refuses. The incompleteness is
therefore at V7 and, by the equivalences just listed, only there.

\emph{Step 6 in the general format ((iv)).} V1 again enumerates declared
tables and compares declared integers and labels: per factor,
$|\SecSet^{f}|$ identifier and $|\SecSet^{f}|$ dimension comparisons, one
block-kind and one length comparison per supplied family against that
factor's map --- \emph{Alice's $\tau_A$, her measurement operators and each
$\Theta_k$ of an entry declared general against the $A$ factor's map among
them}, the same four comparisons each --- plus the comparisons of $\idxm$'s
components with the factors' orders and with $\mathbb{T}$'s. V2, V3 and V6
take the decisions \eqref{eq:predcount-general} counts: $2$ per Alice sector
for V2 --- Hermiticity, which serves both D(ii)'s named condition and the
Hermiticity half of its positivity condition, and semidefiniteness --- and one
trace decision, with one non-zero comparison per sector in ensemble form, all
of these V2 decisions and the marginal trace-summation decision being present
only in trusted mode; when the slot holds $\nomarg$, V2 returns without taking
any of them. V3 has $2|\Xset|+1$ per Alice sector for its Alice half, the
$+1$ being the completeness identity in POVM form and the marginal identity
in ensemble form; $|\SecSet^{f}|(2|\Vset^{f}|+1)$ per measured factor;
$|\SecSet^{f_k}|$ Hermiticity decisions per entry; and
$m_{\text{gen}}|\mathbb{T}|$ for V6. Each is finite and each is taken on a
matrix of one of the four sizes named in the statement, a semidefiniteness on
size $n$ costing $O(n^3)$ exact rational operations as before. V4, V5 and V8
enumerate $|\Xset||\Vset|$ pairs, respectively $(|\Xset||\Vset|)^2$ ordered
pairs, of \emph{tabulated} values, and $a$'s table is finite because each
$a_m$'s is. V7 is Theorem~\ref{thm:feasibility}\,(ii) in its general-format reading,
which is per joint block and in exact arithmetic on both of its branches, for
the reasons given there.
\end{proof}
